\chapter{Cơ Sở Lý Thuyết}

\section{Biểu diễn mạng đường bằng đồ thị}

Mạng đường được mô hình hóa bởi đa đồ thị có hướng có trọng số
$G=(V,E,w)$. Mỗi $v\in V$ có tọa độ $(x_v,y_v)$. Một cạnh
$e=(u,v,k)\in E$ là một đoạn đường từ $u$ tới $v$; chỉ số $k$ cho phép tồn tại
nhiều cạnh song song. Trọng số $w(e)>0$ có thể là chiều dài theo mét hoặc thời
gian ước tính theo giây. Việc dùng graph có hướng giữ được đường một chiều và sự
khác nhau giữa hai hướng di chuyển.

\begin{formulabox}[title={Chi phí của một đường đi}]
Với đường đi $P=(v_0,e_1,v_1,\ldots,e_m,v_m)$, chi phí là
\[
  C(P)=\sum_{i=1}^{m} w(e_i).
\]
Bài toán đường đi ngắn nhất tìm $P^*\in\mathcal{P}(s,t)$ sao cho
\[
  P^*=\argmin_{P\in\mathcal{P}(s,t)} C(P).
\]
\end{formulabox}

Dữ liệu đường được lấy từ OpenStreetMap và xử lý bằng OSMnx
\cite{openstreetmap,boeing2017}. Graph định tuyến đóng gói hiện tại có 350.930
nút, 864.982 cạnh và hộp bao địa lý $[105{,}2;20{,}5;106{,}1;21{,}45]$ theo thứ
tự kinh độ tây, vĩ độ nam, kinh độ đông, vĩ độ bắc.

\section{Thuật toán Dijkstra}

Dijkstra giải bài toán đường đi ngắn nhất nguồn đơn khi mọi trọng số cạnh không
âm \cite{dijkstra1959,cormen2009}. Thuật toán duy trì khoảng cách tạm thời $d[v]$, nút trước
$\pi[v]$ và một hàng đợi ưu tiên. Khi nút có khoảng cách nhỏ nhất được lấy ra và
khoảng cách đó chưa lỗi thời, giá trị của nó đã được chốt.

\begin{formulabox}[title={Phép nới lỏng cạnh}]
Với cạnh $(u,v)$, nếu
\[
  d[u]+w(u,v)<d[v],
\]
thì cập nhật $d[v]\leftarrow d[u]+w(u,v)$ và $\pi[v]\leftarrow u$.
\end{formulabox}

Trong đồ án, Dijkstra nhận một tập đích và có thể dừng khi tất cả các đích đã
được chốt. Với heap nhị phân, độ phức tạp xấp xỉ
$O((|V|+|E|)\log |V|)$ cho mỗi nguồn; bộ nhớ $O(|V|)$. Hàng đợi Fibonacci hạ cận
lý thuyết xuống $O(|E|+|V|\log|V|)$ nhưng hằng số lớn hơn nên không được dùng ở
đây \cite{fredman1987}. Để lập lộ trình qua
$n$ điểm quan tâm, thuật toán được chạy từ từng điểm, tạo ma trận
$D\in\mathbb{R}^{n\times n}$, trong đó $D_{ij}$ là chi phí từ điểm $i$ tới
điểm $j$. Do graph có hướng, nói chung $D_{ij}\neq D_{ji}$.

Điều kiện quan trọng là $w(e)>0$ và các điểm cần thiết phải nối được theo hướng
di chuyển. Nếu không có đường, giá trị được biểu diễn bằng $+\infty$ và tầng tối
ưu thứ tự phải báo lỗi thay vì tạo một tuyến không hợp lệ.

\section{Bài toán thứ tự ghé nhiều điểm}

Sau khi có ma trận $D$, bài toán trở thành một biến thể TSP bất đối xứng. Nút 0
là điểm xuất phát; các nút còn lại là điểm cần ghé. Tùy chế độ, tuyến có thể kết
thúc tự do, kết thúc tại một điểm cố định hoặc quay lại 0.

\subsection{Quy hoạch động Held--Karp}

Held--Karp dùng trạng thái $(S,j)$, trong đó $S$ là tập điểm đã ghé và $j\in S$
là điểm cuối hiện tại \cite{heldkarp1962}. Gọi $F(S,j)$ là chi phí nhỏ nhất để
đi từ 0, ghé đúng các điểm trong $S$ và kết thúc tại $j$.

\begin{formulabox}[title={Truy hồi Held--Karp}]
\[
F(\{j\},j)=D_{0j},
\]
\[
F(S,j)=\min_{i\in S\setminus\{j\}}
  \left[F(S\setminus\{j\},i)+D_{ij}\right].
\]
Với tuyến mở, chọn $\min_j F(U,j)$. Với chu trình, chọn
$\min_j[F(U,j)+D_{j0}]$. Điểm cuối cố định $t$ dùng trực tiếp $F(U,t)$.
\end{formulabox}

Thuật toán có độ phức tạp thời gian $O(n^2 2^n)$ và bộ nhớ $O(n2^n)$
\cite{bellman1962,cormen2009}. Giới hạn
10 điểm ghé của giao diện giữ số trạng thái ở mức phù hợp cho tương tác. Mảng
nút trước được lưu cùng chi phí để khôi phục thứ tự tối ưu, sau đó từng cặp liên
tiếp được thay bằng đường đi Dijkstra tương ứng.

\subsection{Láng giềng gần nhất và 2-opt}
\label{sec:nn-2opt}

Láng giềng gần nhất chọn điểm chưa ghé có chi phí trực tiếp nhỏ nhất
\cite{rosenkrantz1977}. Chi phí
tính toán là $O(n^2)$ nhưng kết quả phụ thuộc mạnh vào lựa chọn cục bộ. 2-opt \cite{croes1958} bắt
đầu từ một tuyến và đảo thứ tự một đoạn nếu phép đổi làm giảm tổng chi phí. Với
ma trận bất đối xứng, hệ thống đánh giá lại toàn bộ chi phí ứng viên thay vì sử
dụng công thức delta dành riêng cho TSP đối xứng.

Khoảng cách thực đo của hai heuristic so với nghiệm chính xác được trình bày ở
mục~\ref{sec:benchmark-tsp}.

Hai heuristic được dùng làm baseline và phương án so sánh, không thay thế
Held--Karp trong miền tối đa 10 điểm. Không có bảo đảm chúng tìm được tối ưu toàn
cục.

\section{Pursuit--evasion trên đồ thị}

Pursuit--evasion trên graph là một họ bài toán trong đó bên truy tìm và bên lẩn
trốn di chuyển trên các nút/cạnh theo luật xác định \cite{parsons1978,chung2011}. Biến thể
của đồ án là rời rạc theo lượt, có lối thoát, ngân sách khoảng cách và quan sát
không đầy đủ.

Trạng thái đầy đủ ở lượt $t$ gồm vị trí thật $r_t$ của người trốn, vị trí các đội
$H_t=(h_t^1,h_t^2)$, tập lối thoát $X$, tập niềm tin $B_t$ và lịch sử cần thiết.
Chính sách máy không nhận toàn bộ trạng thái; đầu vào của nó là quan sát
$O_t=(t,H_t,X,B_t,\ell,\tau,\rho)$, trong đó $\ell$ là vị trí được thấy gần nhất,
$\tau$ là lượt thấy gần nhất và $\rho$ là bộ luật.

\begin{formulabox}[title={Lan truyền tập niềm tin}]
Với ngân sách người trốn $b_r$ và khoảng cách ngắn nhất $d$, tập niềm tin trước
khi áp dụng thông tin mới là
\[
  \widehat{B}_{t+1}=\{v\in V\mid \exists u\in B_t:\ d(u,v)\le b_r\}.
\]
Nếu lượt $t+1$ là lượt tiết lộ thì $B_{t+1}=\{r_{t+1}\}$. Nếu không, các nút
đang bị đội truy tìm chiếm và đã được loại trừ bởi sự sống sót của người trốn sẽ
được bỏ khỏi $\widehat{B}_{t+1}$.
\end{formulabox}

Tập niềm tin là biểu diễn chính xác theo luật chuyển trạng thái đã cài đặt, nhưng
không phải phân phối xác suất. Mọi vị trí khả dĩ được coi như nhau; vì vậy chính
sách coverage là heuristic dựa trên khoảng cách, không phải suy luận Bayes. Dạng tổng quát của bài toán ra quyết định dưới quan sát
một phần là POMDP, trong đó trạng thái niềm tin là một phân phối xác suất chứ
không phải một tập \cite{kaelbling1998}.

\section{Ba chiến thuật truy tìm}

Gọi $b_h$ là ngân sách mỗi đội. Điểm bao phủ của vị trí ứng viên $q$ đối với tập
$S$ được định nghĩa:

\begin{formulabox}[title={Điểm bao phủ nghịch đảo khoảng cách}]
\[
  \operatorname{cov}(q,S)=\frac{1}{|S|}\sum_{p\in S}
  \frac{1}{1+d(q,p)/b_h}.
\]
\end{formulabox}

\textbf{Đuổi dấu vết} đưa mỗi đội theo đường ngắn nhất về vị trí được thấy gần
nhất. \textbf{Tìm kiếm độc lập} để mỗi đội chọn riêng vị trí khả đạt có điểm
bao phủ lớn nhất. \textbf{Phối hợp chặn lối thoát} xét tích Descartes của tối đa
tám ứng viên tốt nhất cho mỗi đội, bổ sung các bước hướng lối thoát, rồi tối đa
hóa:

\begin{formulabox}[title={Điểm phối hợp}]
\[
J(M)=\frac{1}{|B|}\sum_{p\in B}\max_{m\in M}
\frac{1}{1+d(m,p)/b_h}
+0{,}35\cdot\frac{1}{|X|}\sum_{x\in X}\max_{m\in M}
\frac{1}{1+d(m,x)/b_h}.
\]
\end{formulabox}

Hệ số 0,35 là tham số thiết kế của phiên bản hiện tại, chưa được tối ưu bằng một
tập validation độc lập. Việc cắt còn tám ứng viên giữ thời gian phản hồi nhỏ,
nhưng cũng có thể loại một nước đi tốt; đây là một giới hạn có chủ ý của heuristic.

\section{Bất biến và thứ tự ưu tiên kết thúc}

Các bất biến được dùng để kiểm thử gồm: người trốn luôn thuộc tập niềm tin khi
ván đang diễn ra; mọi bước đi nằm trong ngân sách và đúng chiều cạnh; chính sách
máy không nhận vị trí ẩn; tập niềm tin không rỗng trong trạng thái hợp lệ; và
trạng thái kết thúc không được đi tiếp.

